\subsection{RGB Image}

An RGB image represents colour using three components: Red, Green, and
Blue. A pixel can therefore be represented by the tuple
\begin{align}
P_{ij} &= (R_{ij},G_{ij},B_{ij})\label{eq:rgb_pixel}
\end{align}

For a conventional 8-bit RGB image, each component is represented using
8 bits. Hence, one pixel requires
\begin{align}
8+8+8 &= 24\label{eq:rgb_bits}
\end{align}
bits of colour information.

For an image containing $M\times N$ pixels, the image may be viewed as
three matrices,
\begin{align}
R&\in\mathbb{R}^{M\times N}\label{eq:r_matrix}\\
G&\in\mathbb{R}^{M\times N}\label{eq:g_matrix}\\
B&\in\mathbb{R}^{M\times N}\label{eq:b_matrix}
\end{align}

The conversion to black and white is therefore a reduction from three
colour values per pixel to one intensity value per pixel.
